\section{Tractable Special Cases}\label{sec:tractable}

We distinguish the encodings of Section~\ref{sec:encoding}. The tractability results below state the model assumption explicitly. All results are formalized in \texttt{DecisionQuotient/Tractability/} and verified in Lean~4 with no \texttt{sorry} placeholders.

\begin{theorem}[Tractable Subcases]\label{thm:tractable}
SUFFICIENCY-CHECK is polynomial-time solvable in the following cases:
\begin{enumerate}
\item \textbf{Bounded actions (explicit-state encoding):} The input contains the full utility table over $A \times S$ with $|A| \leq k$ for constant $k$. Complexity: $O(|S|^2 \cdot k^2)$.
\item \textbf{Separable utility (rank 1):} $U(a, s) = f(a) + g(s)$. Complexity: $O(1)$ (trivially sufficient).
\item \textbf{Low tensor rank:} The utility admits a rank-$R$ decomposition $U(a,s) = \sum_{r=1}^R w_r \cdot f_r(a) \cdot \prod_i g_{ri}(s_i)$. Complexity: $O(|A| \cdot R \cdot n)$.
\item \textbf{Tree-structured dependencies:} Coordinates are ordered such that $s_i$ depends only on $(s_1, \ldots, s_{i-1})$. Complexity: $O(n \cdot |A| \cdot k_{\max})$ where $k_{\max} = \max_i |X_i|$.
\item \textbf{Bounded treewidth:} The interaction graph has treewidth $\leq w$ and utility factors over edges. Complexity: $O(n \cdot k^{w+1})$ via standard CSP algorithms.
\item \textbf{Coordinate symmetry:} Utility is invariant under coordinate permutations. Complexity: $O\bigl(\binom{d+k-1}{k-1}^2\bigr)$ orbit-type pairs.
\end{enumerate}
\end{theorem}

The following subsections provide formal definitions, reduction theorems, and prose descriptions for each tractable subcase. For each result, we cite standard algorithms where applicable and prove only the paper-specific reductions in Lean.

%% ============================================================
\subsection{Bounded Actions}\label{sec:tract-bounded}

When the action set is bounded, brute-force enumeration becomes polynomial.

\begin{definition}[Bounded Action Problem]
A decision problem has \emph{bounded actions} with parameter $k$ if $|A| \leq k$ for some constant $k$ independent of the state space size.
\end{definition}

\begin{theorem}[Bounded Actions Tractability]\label{thm:bounded-actions}
If $|A| \leq k$ for constant $k$, then SUFFICIENCY-CHECK runs in $O(|S|^2 \cdot k^2)$ time.
\end{theorem}

\begin{proof}
Given the full table of $U(a,s)$:
\begin{enumerate}
\item Compute $\Opt(s)$ for all $s \in S$ in $O(|S| \cdot k)$ time.
\item For each pair $(s, s')$ with $s_I = s'_I$, compare $\Opt(s)$ and $\Opt(s')$ in $O(k^2)$ time.
\item Total: $O(|S|^2)$ pairs $\times$ $O(k^2)$ comparisons = $O(|S|^2 \cdot k^2)$.
\end{enumerate}
When $k$ is constant, this is polynomial in $|S|$.
\end{proof}

\noindent\textbf{Lean formalization:} \texttt{DecisionQuotient/Tractability/BoundedActions.lean}
\begin{itemize}
\item \LH{L10}: Correctness of the algorithm.
\item \LH{L2}: Complexity bound $|S|^2 \cdot (1 + k^2)$.
\end{itemize}

%% ============================================================
\subsection{Separable Utility (Rank 1)}\label{sec:tract-separable}

Separable utility decouples actions from states, making sufficiency trivial.

\begin{definition}[Separable Utility]\label{def:separable}
A utility function is \emph{separable} if there exist functions $f : A \to \mathbb{R}$ and $g : S \to \mathbb{R}$ such that $U(a,s) = f(a) + g(s)$ for all $(a,s) \in A \times S$.
\end{definition}

\begin{theorem}[Separable Tractability]\label{thm:separable}
If utility is separable, then $I = \emptyset$ is always sufficient.
\end{theorem}

\begin{proof}
If $U(a, s) = f(a) + g(s)$, then:
\[
\Opt(s) = \arg\max_{a \in A} [f(a) + g(s)] = \arg\max_{a \in A} f(a)
\]
The optimal action depends only on $f$, not on $s$. Hence $\Opt(s) = \Opt(s')$ for all $s, s'$, making any $I$ (including $\emptyset$) sufficient.
\end{proof}

\noindent\textbf{Lean formalization:} \texttt{DecisionQuotient/Tractability/SeparableUtility.lean}
\begin{itemize}
\item \LH{SeparableUtility}: Structure encoding the separability condition.
\item \LH{L8}: Optimal actions are state-independent.
\item \LH{L8}: The empty set is sufficient.
\end{itemize}


%% ============================================================
\subsection{Low Tensor Rank}\label{sec:tract-tensor}

Utility functions with low tensor rank admit efficient factored computation, generalizing separable utility.

\begin{definition}[Tensor Rank Decomposition]\label{def:tensor-rank}
A utility function $U : A \times ((i : \mathsf{Fin}\ n) \to X_i) \to \mathbb{R}$ has \emph{tensor rank} $R$ if there exist weights $w_1, \ldots, w_R \in \mathbb{R}$, action factors $f_r : A \to \mathbb{R}$, and coordinate factors $g_{ri} : X_i \to \mathbb{R}$ such that:
\[
U(a, s) = \sum_{r=1}^R w_r \cdot f_r(a) \cdot \prod_{i=1}^n g_{ri}(s_i)
\]
\end{definition}

\begin{remark}
Separable utility (Definition~\ref{def:separable}) is the special case $R = 1$ with $w_1 = 1$, $f_1(a) = f(a)$, and $\prod_i g_{1i}(s_i) = g(s)$.
\end{remark}

\begin{theorem}[Low Tensor Rank Tractability]\label{thm:tensor-rank}
If utility has tensor rank $R$, then SUFFICIENCY-CHECK runs in $O(|A| \cdot R \cdot n)$ time per state.
\end{theorem}

\begin{proof}
The factored representation enables efficient computation:
\begin{enumerate}
\item For each rank component $r$, the contribution $w_r \cdot f_r(a) \cdot \prod_i g_{ri}(s_i)$ factors over coordinates.
\item Two states $s, s'$ agreeing on $I$ have $g_{ri}(s_i) = g_{ri}(s'_i)$ for $i \in I$.
\item The partial product $\prod_{i \in I} g_{ri}(s_i)$ is identical for agreeing states.
\item This reduces sufficiency checking to comparing factored sums over orbit representatives.
\end{enumerate}
\end{proof}

\noindent\textbf{Lean formalization:} \texttt{DecisionQuotient/Tractability/SeparableUtility.lean}
\begin{itemize}
\item \LH{TensorRankDecomposition}: Structure encoding rank-$R$ decomposition.
\item \LH{L14}: Axiom citing tensor network algorithms~\cite{kolda2009tensor,cichocki2016tensor}.
\item \LH{L4}: Reduction theorem for factored computation.
\item \LH{L6}: Tractability via factored sufficiency check.
\end{itemize}

%% ============================================================
\subsection{Tree-Structured Dependencies}\label{sec:tract-tree}

When coordinate dependencies form a tree, dynamic programming yields polynomial-time sufficiency checking.

\begin{definition}[Tree-Structured Dependencies]\label{def:tree-struct}
A decision problem has \emph{tree-structured dependencies} if there exists an ordering of coordinates $(1, \ldots, n)$ such that the utility decomposes as:
\[
U(a, s) = \sum_{i=1}^n u_i(a, s_i, s_{\mathrm{parent}(i)})
\]
where $\mathrm{parent}(i) < i$ for all $i > 1$, and the root term depends only on $(a, s_1)$.
\end{definition}

\begin{theorem}[Tree Structure Tractability]\label{thm:tree-struct}
If dependencies are tree-structured with explicit local tables, SUFFICIENCY-CHECK runs in $O(n \cdot |A| \cdot k_{\max})$ time where $k_{\max} = \max_i |X_i|$.
\end{theorem}

\begin{proof}
Dynamic programming on the tree order:
\begin{enumerate}
\item Process coordinates in order $i = n, n-1, \ldots, 1$.
\item For each node $i$ and each value of its parent coordinate, compute the set of actions optimal for some subtree assignment.
\item Combine child summaries with local tables in $O(|A| \cdot k_{\max})$ per node.
\item Total: $O(n \cdot |A| \cdot k_{\max})$.
\end{enumerate}
A coordinate is relevant iff varying its value changes the optimal action set.
\end{proof}

\noindent\textbf{Lean formalization:} \texttt{DecisionQuotient/Tractability/TreeStructure.lean}
\begin{itemize}
\item \LH{TreeStructured}: Predicate encoding tree dependency structure.
\item \LH{L16}: Axiom citing standard DP on trees.
\end{itemize}

%% ============================================================
\subsection{Bounded Treewidth}\label{sec:tract-treewidth}

The tree-structured case generalizes to bounded treewidth interaction graphs via CSP reduction.

\begin{definition}[Interaction Graph]\label{def:interaction-graph}
Given a pairwise utility decomposition, the \emph{interaction graph} $G = (V, E)$ has:
\begin{itemize}
\item Vertices $V = \{1, \ldots, n\}$ (one per coordinate)
\item Edges $E = \{(i, j) : i \text{ and } j \text{ interact in the utility}\}$
\end{itemize}
\end{definition}

\begin{definition}[Pairwise Utility]\label{def:pairwise}
A utility function is \emph{pairwise} if it decomposes as:
\[
U(a, s) = \sum_{i=1}^n u_i(a, s_i) + \sum_{(i,j) \in E} u_{ij}(a, s_i, s_j)
\]
with unary terms $u_i$ and binary terms $u_{ij}$ given explicitly.
\end{definition}

\begin{theorem}[Bounded Treewidth Tractability]\label{thm:treewidth}
If the interaction graph has treewidth $\leq w$ and utility is pairwise with explicit factors, then SUFFICIENCY-CHECK runs in $O(n \cdot k^{w+1})$ time where $k = \max_i |X_i|$.
\end{theorem}

\begin{proof}
We reduce SUFFICIENCY-CHECK to constraint satisfaction on the interaction graph:
\begin{enumerate}
\item \textbf{Reduction:} For each action $a \in A$, checking whether $\Opt(s) = \Opt(s')$ for states agreeing on $I$ reduces to a CSP on $G$ with variables $\{s_i : i \notin I\}$.
\item \textbf{Standard algorithm:} CSP on graphs of treewidth $w$ is solvable in $O(n \cdot k^{w+1})$ via tree decomposition~\cite{bodlaender1993tourist,freuder1990complexity}.
\end{enumerate}
The reduction is polynomial; the algorithm is standard.
\end{proof}

\noindent\textbf{Lean formalization:} \texttt{DecisionQuotient/Tractability/TreeStructure.lean}
\begin{itemize}
\item \LH{InteractionGraph}: Definition of the interaction graph.
\item \LH{PairwiseUtility}: Structure encoding pairwise decomposition.
\item \LH{L16}: Axiom for treewidth (citing Bodlaender~\cite{bodlaender1993tourist}).
\item \LH{L4}: Axiom citing CSP algorithms.
\item \LH{L12}: Paper-specific reduction theorem.
\item \LH{L2}: Tractability statement combining reduction and algorithm.
\end{itemize}


%% ============================================================
\subsection{Coordinate Symmetry}\label{sec:tract-symmetry}

When utility is invariant under coordinate permutations, the effective state space collapses to orbit types.

\begin{definition}[Dimensional State Space]\label{def:dimensional}
A \emph{dimensional state space} with alphabet size $k$ and dimension $d$ is $S = \{0, \ldots, k-1\}^d$, the set of $d$-tuples over a $k$-element alphabet.
\end{definition}

\begin{definition}[Symmetric Utility]\label{def:symmetric}
A utility function on a dimensional state space is \emph{symmetric} if it is invariant under coordinate permutations: for all permutations $\sigma \in \mathfrak{S}_d$ and states $s \in S$,
\[
U(a, s) = U(a, \sigma \cdot s) \quad \text{where } (\sigma \cdot s)_i = s_{\sigma^{-1}(i)}
\]
\end{definition}

\begin{definition}[Orbit Type]\label{def:orbit-type}
The \emph{orbit type} of a state $s \in \{0, \ldots, k-1\}^d$ is the multiset of its coordinate values:
\[
\text{orbitType}(s) = \{\!\{s_1, s_2, \ldots, s_d\}\!\}
\]
Two states have the same orbit type iff they lie in the same orbit under coordinate permutation.
\end{definition}

\begin{theorem}[Orbit Type Characterization]\label{thm:orbit-type}
For dimensional state spaces, two states $s, s'$ have equal orbit types if and only if there exists a coordinate permutation $\sigma$ such that $\sigma \cdot s = s'$.
\end{theorem}

\begin{proof}
\textbf{If:} Permuting coordinates preserves the multiset of values.
\textbf{Only if:} If $s$ and $s'$ have the same multiset of values, we can construct a permutation mapping each occurrence in $s$ to the corresponding occurrence in $s'$. The Lean proof constructs this permutation explicitly via a fiber bundle argument.
\end{proof}

\begin{theorem}[Symmetric Sufficiency Reduction]\label{thm:symmetric-reduction}
Under symmetric utility, optimal actions depend only on orbit type:
\[
\text{orbitType}(s) = \text{orbitType}(s') \implies \Opt(s) = \Opt(s')
\]
Thus SUFFICIENCY-CHECK reduces to checking pairs with \emph{different} orbit types.
\end{theorem}

\begin{proof}
If $\text{orbitType}(s) = \text{orbitType}(s')$, then by Theorem~\ref{thm:orbit-type} there exists $\sigma$ with $\sigma \cdot s = s'$. By symmetry, $U(a, s) = U(a, s')$ for all $a$, so $\Opt(s) = \Opt(s')$.

For sufficiency, we need only check pairs $(s, s')$ agreeing on $I$ with \emph{different} orbit types; same-orbit pairs are automatically equal.
\end{proof}

\begin{theorem}[Symmetric Complexity Bound]\label{thm:symmetric-complexity}
The number of orbit types is bounded by the stars-and-bars count:
\[
|\text{OrbitTypes}| = \binom{d + k - 1}{k - 1}
\]
Thus SUFFICIENCY-CHECK requires at most $\binom{d + k - 1}{k - 1}^2$ orbit-type pair checks, which is polynomial in $d$ for fixed $k$.
\end{theorem}

\begin{proof}
An orbit type is a multiset of $d$ values from $\{0, \ldots, k-1\}$, equivalently a $k$-tuple of non-negative integers summing to $d$. By stars-and-bars, the count is $\binom{d + k - 1}{k - 1}$.

For fixed $k$, this is $O(d^{k-1})$, polynomial in $d$.
\end{proof}

\noindent\textbf{Lean formalization:} \texttt{DecisionQuotient/Tractability/Dimensional.lean}
\begin{itemize}
\item \LH{DimensionalStateSpace}: Definition of $k$-ary $d$-dimensional state space.
\item \LH{SymmetricUtility}: Predicate for coordinate-permutation invariance.
\item \LH{orbitType}: Function computing the orbit type (histogram).
\item \LH{L6}: \textbf{(Core result)} Orbit types equal iff permutation exists.
\item \LH{L12}: Optimal actions constant on orbits.
\item \LH{L10}: Reduction to cross-orbit pairs.
\item \LH{L14}: $O\bigl(\binom{d+k-1}{k-1}^2\bigr)$ bound.
\end{itemize}

%% ============================================================
\subsection{Practical Implications}\label{sec:tract-practical}

The six tractable cases correspond to common modeling scenarios:

\begin{itemize}
\item \textbf{Bounded actions:} Small action sets (e.g., binary decisions, few alternatives).
\item \textbf{Separable utility:} Additive cost models, linear utility functions.
\item \textbf{Low tensor rank:} Factored preference models, low-rank approximations.
\item \textbf{Tree structure:} Hierarchical decision processes, causal models with tree structure.
\item \textbf{Bounded treewidth:} Spatial models, graphical models with limited interaction.
\item \textbf{Coordinate symmetry:} Exchangeable features, order-invariant utilities.
\end{itemize}

For problems given in the succinct encoding without these structural restrictions, the hardness results of Section~\ref{sec:hardness} apply, justifying heuristic approaches.

%% ============================================================
\subsection{Verification Philosophy}\label{sec:tract-verification}

The Lean formalization follows a consistent pattern for each tractable case:

\begin{enumerate}
\item \textbf{Define the structural restriction} as a Lean structure or predicate.
\item \textbf{State axioms for standard algorithms} (DP on trees, CSP on bounded treewidth, tensor contraction) with citations to the literature.
\item \textbf{Prove the paper-specific reduction} showing SUFFICIENCY-CHECK satisfies the preconditions of the standard algorithm.
\item \textbf{Combine} to obtain the tractability statement.
\end{enumerate}

This separation ensures that:
\begin{itemize}
\item Standard results are cited, not re-proved (avoiding redundant mechanization).
\item Paper-specific claims are fully verified (the reductions are where errors hide).
\item The formalization is maintainable (algorithm updates don't require re-verification).
\end{itemize}

The philosophy: \emph{mechanize what requires mechanical verification for rational belief; cite what is standard}.
